\documentclass[11pt]{article}
\usepackage[T1]{fontenc}
\usepackage{lmodern}
\usepackage[margin=1in]{geometry}
\usepackage{amsmath,amssymb,amsthm,mathtools}
\usepackage{microtype,booktabs,array,enumitem}
\usepackage[dvipsnames]{xcolor}
\usepackage{url}
\usepackage[colorlinks=true,linkcolor=MidnightBlue,citecolor=MidnightBlue,urlcolor=MidnightBlue]{hyperref}
\hypersetup{pdftitle={Finite-horizon Weil coercivity and shifted-zeta contraction},pdfauthor={Edward Baker}}
\numberwithin{equation}{section}
\newtheorem{theorem}{Theorem}[section]
\newtheorem{proposition}[theorem]{Proposition}
\newtheorem{lemma}[theorem]{Lemma}
\newtheorem{corollary}[theorem]{Corollary}
\theoremstyle{definition}
\newtheorem{definition}[theorem]{Definition}
\theoremstyle{remark}
\newtheorem{remark}[theorem]{Remark}
\newcommand{\HH}{\mathcal H}
\newcommand{\DD}{\mathcal D}
\newcommand{\RR}{\mathbb R}
\newcommand{\CC}{\mathbb C}
\newcommand{\Rea}{\operatorname{Re}}
\newcommand{\supp}{\operatorname{supp}}
\newcommand{\tr}{\operatorname{tr}}
\newcommand{\norm}[1]{\lVert#1\rVert}
\newcommand{\ip}[2]{\langle#1,#2\rangle}
\newcommand{\gam}{\Gamma_\infty}
\newcommand{\Lstar}{L_*}
\newcommand{\mstar}{m_*}
\newcommand{\hstar}{h_*}
\setlength{\emergencystretch}{2em}
\setlist{itemsep=3pt,topsep=5pt}
\title{Finite-horizon Weil coercivity\protect\\ and shifted-zeta contraction\protect\\[5pt]
\large Full-output parity certificates through total horizon $\log 7$}
\author{Edward Baker}
\date{Review draft, version 0.2\quad---\quad 9 September 2026}

\begin{document}
\maketitle
\begin{abstract}
We study the central Weil form associated with the causal finite-horizon
transfer of Suzuki's shifted zeta ratio. A computer-assisted calculation proves
$Q_{0,L}\succeq10^{-26}I$ for every total horizon $0<L\le9/5$, and
$Q_{0,L}\succeq10^{-34}I$ for $0<L\le\log7$.
These horizons lie beyond $\log6$: the generator contains the prime-power
delays $2,3,4,5$, while the transfer includes the first mixed-prime composition.
The proof combines an analytic lower bound on the infinite Legendre tail with
full-output Gram matrices retaining the exact reflection-parity coupling.
Both resulting $64\times64$ Schur inequalities are verified by ball arithmetic,
with every finite construction operation enclosed and an explicit analytic
profile remainder subtracted. An even-in-shift generator identity gives
$\norm{V_{\omega,L}}\le\exp(-5\times10^{-27}\omega)$ for
$0<\omega\le3\times10^{-14}$ and $L\le9/5$, together with quantitative
reflection-positivity and cumulative-storage bounds. At the extended horizon,
the same argument gives $\norm{V_{\omega,L}}\le\exp(-5\times10^{-35}\omega)$
for $0<\omega\le3\times10^{-18}$ and $L\le\log7$.
An independent implementation reproduces both certificates. Exact storage
and nested Schur identities specify what an induction would need to retain.
The argument does not use zeta zeros or an assumed positive spectral measure. It establishes a
finite-depth result, not a zero-free region or an all-depth positivity theorem.
\end{abstract}

\noindent\textbf{Keywords.} Weil positivity; Riemann zeta function; causal
convolution; Hankel operators; Schur complement; validated computation.

\smallskip
\noindent\textbf{Review status.} The analytic argument and executable
certificate are supplied for review. The finite calculation and replay have
passed; independent human verification and proof-assistant formalization have
not been completed. The numerical results are isolated in
Theorems~\ref{thm:computation} and~\ref{thm:log7}.

\section{Introduction and main result}

Weil's positivity criterion expresses the Riemann hypothesis as positivity of
an explicit quadratic form on all compactly supported test functions. Fixing
a bounded support produces a separate, finite-horizon question: is the
restricted form coercive, and can an explicit positive lower bound be proved
without information about zeta zeros? Such a result controls only a bounded
part of the criterion, but provides a test of methods intended to extend to
larger arithmetic depth.

This paper addresses that question using the causal transfer associated with
the shifted quotient
\begin{equation}\label{eq:theta}
 \Theta_\omega(z)=
 \frac{\xi(\tfrac12-\omega-iz)}{\xi(\tfrac12+\omega-iz)},\qquad
 \xi(s)=\tfrac12s(s-1)\pi^{-s/2}\Gamma(s/2)\zeta(s).
\end{equation}
Suzuki introduced this family and its arithmetic Hankel kernel
\cite{Suzuki2016}. Our statements concern the finite causal convolution
defined in Section~\ref{sec:operators}; no global innerness assumption is made.

Throughout the paper, $L$ is the \emph{total length} of the input interval.
Its centered version is $(-L/2,L/2)$, and its autocorrelation is supported
in $[-L,L]$. The corresponding multiplicative cutoff is $a=e^{L/2}$.
This convention is important when comparing arithmetic thresholds or
compact-window bounds.

\begin{theorem}[Finite-horizon coercivity and contraction]\label{thm:main}
Let
\[
 \Lstar=\frac95,\qquad \mstar=10^{-26},\qquad \hstar=3\times10^{-14}.
\]
For every $0<L\le\Lstar$, the central form defined in
\eqref{eq:weil} satisfies
\begin{equation}\label{eq:maincentral}
 Q_{0,L}[f]\ge\mstar\norm f_2^2\qquad(f\in\DD_{\log,L}).
\end{equation}
For the transfer $V_{\omega,L}$ of Definition~\ref{def:transfer},
\begin{equation}\label{eq:maintransfer}
 \norm{V_{\omega,L}}\le e^{-5\times10^{-27}\omega}<1,
 \qquad 0<\omega\le\hstar.
\end{equation}
Consequently $H_{\omega,L}=V_{\omega,L}R_L$ is self-adjoint and
\begin{equation}\label{eq:mainreflection}
 I\pm H_{\omega,L}\succeq
 (1-e^{-5\times10^{-27}\omega})I.
\end{equation}
The central estimate is computer-assisted; the remaining implications are
analytic.
\end{theorem}

\begin{theorem}[The logarithmic endpoint]\label{thm:endpoint-main}
For every $0<L\le\log7$,
\[
 Q_{0,L}\succeq10^{-34}I,\qquad
 \norm{V_{\omega,L}}\le e^{-5\times10^{-35}\omega}
 \quad(0<\omega\le3\times10^{-18}).
\]
Consequently $I\pm H_{\omega,L}\succeq(1-e^{-5\times10^{-35}\omega})I$.
The sharper bounds of Theorem~\ref{thm:main} remain available for $L\le9/5$.
\end{theorem}

The total horizon $9/5$ is strictly larger than $\log6$. At this horizon
the logarithmic generator has precisely the delays $\log2,\log3,\log4,\log5$.
The coefficient of $n=6$ occurs in the transfer as a product of two prime
factors, although $\Lambda(6)=0$. In the full-output calculation the ordered
mixed overlaps $(2,3)$ and $(3,2)$ are retained explicitly.

The proof has three components. First, differentiating the finite Euler
product gives a generator even in $\omega$ and a bounded perturbation
$Q_{\omega,L}-Q_{0,L}=O_L(\omega^2)$. Second, a fractional-integration estimate
gives a lower bound on the infinite Legendre tail. Third, full-output Grams
control the coupling between that tail and the retained space, separately in
each reflection sector. The finite comparison is a Schur inequality at the
claimed positive level $\mstar$; positivity of a finite compression alone
would not suffice.

\subsection{Relation to earlier work}

Connes and Consani study archimedean Weil positivity through a compressed
scaling action and Sonin spaces \cite{ConnesConsani2020}. Their work motivates
the broader question of combining archimedean and arithmetic terms.
The present construction uses causal shift differentiation and a polynomial
tail estimate.

The recent preprint of Zhu \cite{Zhu2026}, version 2, reports certified
compact-window positivity at total input length $1.6$, using a
frequency-envelope reduction. Its window parameter is a half-width, so its
$L=0.8$ corresponds to total length $1.6$ here. Our central form has the same
geometric normalization after that conversion. We do not use its numerical
certificate. The specific contribution here is the full-output parity
reduction, its fractional tail bound, and an independently assembled
certificates through total lengths $1.8$ and $\log7$, with explicit positive-shift
consequences. We make no claim that finite-matrix approaches to Weil
positivity are themselves new.

The estimates are finite-depth. In particular, they neither establish global
innerness of \eqref{eq:theta} nor prove a new zero-free region. The shift
interval is small, and the storage corollary below does not supply an induction
in arithmetic depth.

\section{Causal transfer and the central form}\label{sec:operators}

\subsection{Finite-horizon definitions}

Let $\HH_L=L^2(0,L;\CC)$, with inner product linear in the second argument,
and define
\[
 R_Lf(x)=f(L-x),\qquad T_df(x)=\mathbf1_{x>d}f(x-d).
\]
In particular, $T_d=0$ for $d\ge L$ and
$T_dT_e=T_{d+e}$. Every adjoint and norm below refers to unweighted
$L^2(0,L)$.

Put $\gam(s)=\tfrac12s(s-1)\pi^{-s/2}\Gamma(s/2)$. For $0<\omega\le1/2$
the archimedean causal convolution has Laplace transfer
\begin{equation}\label{eq:gammaK}
 K^\gamma_\omega(p)=
 \frac{\gam(\tfrac12+p-\omega)}{\gam(\tfrac12+p+\omega)}.
\end{equation}
It may be defined on a sufficiently far right Laplace line, followed by
restriction to $(0,L)$. For an explicit impulse, set
$\alpha=\tfrac12-\omega$, $\beta=\tfrac12+\omega$. Then
\begin{equation}\label{eq:gammafactor}
 K^\gamma_\omega(p)=\pi^\omega
 \frac{(p+\alpha)(p-\beta)}{(p-\alpha)(p+\beta)}
 \frac{\Gamma((p+\alpha)/2)}{\Gamma((p+\beta)/2)}.
\end{equation}
The last ratio is the transform of
$2e^{-\alpha t}(1-e^{-2t})^{\omega-1}/\Gamma(\omega)$.
The rational factor has an impulse consisting of a delta mass and smooth
exponentials (with the usual limiting interpretation at $\alpha=0$).
Thus the full impulse $\kappa_\omega$ is locally integrable, behaves as
$O(t^{\omega-1})$ at the origin, and defines a bounded finite-horizon
convolution $V^\gamma_{\omega,L}$.

\begin{definition}\label{def:transfer}
For $0<\omega\le1/2$, define
\begin{align}
 b_\omega(n)&=n^{\omega-1/2}\prod_{p\mid n}(1-p^{-2\omega}),
 & b_\omega(1)&=1,\label{eq:bcoeff}\\
 B_{\omega,L}&=\sum_{\log n<L}b_\omega(n)T_{\log n},
 &V_{\omega,L}&=B_{\omega,L}V^\gamma_{\omega,L}.\label{eq:Vdef}
\end{align}
Set $V_{0,L}=I$ in the strong-limit sense proved below.
\end{definition}

The sum is finite and contains all integers below $e^L$. This is the
logarithmic-coordinate form of Suzuki's arithmetic kernel. The Euler product
identity for its Dirichlet series also identifies the full Laplace transfer,
on a sufficiently far right line, as
$\xi(\tfrac12+p-\omega)/\xi(\tfrac12+p+\omega)$.
Real causal convolution gives $V_{\omega,L}^*=R_LV_{\omega,L}R_L$;
therefore $H_{\omega,L}=V_{\omega,L}R_L$ is self-adjoint and
$\norm{H_{\omega,L}}=\norm{V_{\omega,L}}$.

\subsection{The central Weil form}

For $f\in\HH_L$, let $f_c(x)=f(x+L/2)$ on $(-L/2,L/2)$ and zero elsewhere.
Write
\[
 \widehat f_c(\tau)=\int_{\RR}f_c(x)e^{-i\tau x}\,dx,\qquad
 C(f)=\int_{\RR}f_c(x)\cosh(x/2)\,dx,\quad
 S(f)=\int_{\RR}f_c(x)\sinh(x/2)\,dx.
\]
The logarithmic form domain is
\begin{equation}\label{eq:domain}
 \DD_{\log,L}=\left\{f\in\HH_L:
 \int_{\RR}\log(2+|\tau|)|\widehat f_c(\tau)|^2\,d\tau<\infty\right\}.
\end{equation}
With $\psi=\Gamma'/\Gamma$, define
\begin{align}
 Q_{0,L}[f]={}&\frac1{2\pi}\int_{\RR}
 \bigl(\Rea\psi(\tfrac14+\tfrac{i\tau}{2})-\log\pi\bigr)
 |\widehat f_c(\tau)|^2\,d\tau\notag\\
 &+2|C(f)|^2-2|S(f)|^2
 -\sum_{\log n<L}\frac{\Lambda(n)}{\sqrt n}
 \ip f{(T_{\log n}+T_{\log n}^*)f}.
 \label{eq:weil}
\end{align}
Here $\Lambda(p^k)=\log p$ and is zero otherwise. The last scalar is real.
The Fourier multiplier differs from $\log(2+|\tau|)$ by a bounded function,
by the digamma asymptotics \cite[\S5.11]{DLMF}. All remaining terms are
bounded. Hence \eqref{eq:weil} is a closed semibounded form on
\eqref{eq:domain}. It is the central geometric Weil form in this normalization.
For example, if $g=f_c*\widetilde f_c$ with
$\widetilde f_c(x)=\overline{f_c(-x)}$, its pole contribution is
$\int g(t)(e^{t/2}+e^{-t/2})\,dt=2|C(f)|^2-2|S(f)|^2$.
Its arithmetic contribution is
$-\sum\Lambda(n)n^{-1/2}(g(\log n)+g(-\log n))$, exactly the delay term
above. This also makes translation invariance explicit for complex tests.
It commutes with reflection and with complex conjugation. The same notation
$Q_{0,L}$ denotes its associated self-adjoint operator where appropriate.

\begin{lemma}[Polynomial core]\label{lem:core}
Polynomials restricted to $(0,L)$ form a core for $Q_{0,L}$.
If $P_N$ is the orthogonal projection onto the polynomials of degree less
than $N$, polynomial combinations in $(I-P_N)\HH_L$ are a form core for
the restriction to that subspace.
\end{lemma}
\begin{proof}
Equip \eqref{eq:domain} with its weighted Fourier norm, equivalent to a
shifted form norm. First shrink the centered zero extension by a dilation
about the origin so that its support lies in a smaller interval. Dilations
tending to the identity converge in this weighted norm: their norms are
uniformly bounded, since the logarithmic weights at dilated frequencies are
comparable, and strong convergence follows first on smooth Fourier data and
then by density. Mollification with support smaller than the resulting
boundary gap converges in the weighted norm by dominated convergence.
Consequently $C_c^\infty(0,L)$ is a form core.

For $h\in C^1[0,L]$, integration by parts in its zero extension gives
\[
 |\widehat h(\tau)|\le L\norm h_\infty,\qquad
 |\widehat h(\tau)|\le
 \frac{2\norm h_\infty+L\norm{h'}_\infty}{|\tau|}\quad(\tau\ne0).
\]
These bounds imply continuity from $C^1[0,L]$ to the logarithmic form
domain. Approximate the derivative of a smooth function uniformly by a
polynomial and integrate, matching one value. This gives polynomial
approximation in $C^1$, hence in form norm. Endpoint jumps of the zero
extension cause no difficulty.

Finally $P_N$ is bounded in form norm: it has finite rank, its range consists
of form-domain polynomials, and its coefficient functionals are bounded in
$L^2$. Apply $I-P_N$ to the polynomial approximants of a vector orthogonal
to the retained space.
\end{proof}

\section{The prime-delay generator and continuation}\label{sec:generator}

\subsection{Exact logarithmic differentiation}

Write $\ell(p)=\log\gam(\tfrac12+p)$ on a far right half-plane. The
archimedean logarithmic generator has symbol
\begin{align}
 a^\gamma_\omega(p)&=\ell'(p-\omega)+\ell'(p+\omega)\notag\\
 &=\frac1{p+\alpha}+\frac1{p-\alpha}
 +\frac1{p+\beta}+\frac1{p-\beta}-\log\pi\notag\\
 &\hspace{18mm}+\frac12\psi((p+\alpha)/2)+\frac12\psi((p+\beta)/2).
 \label{eq:agamma}
\end{align}
It defines a causal logarithmic operator $A^\gamma_{\omega,L}$ on its natural
core. Its symmetric form at zero is the gamma part of \eqref{eq:weil}.
For example, the two central rational terms give the symmetric kernel
$2\cosh((x-y)/2)$, whose form is $2|C|^2-2|S|^2$; the digamma term gives
the multiplier in \eqref{eq:weil}.

\begin{proposition}[Prime-power generator]\label{prop:generator}
On the causal generator core,
\begin{align}
 \partial_\omega V_{\omega,L}&=-A_{\omega,L}V_{\omega,L},\label{eq:derivative}\\
 A_{\omega,L}&=A^\gamma_{\omega,L}
 -2\sum_{\log n<L}\frac{\Lambda(n)}{\sqrt n}
 \cosh(\omega\log n)T_{\log n}.\label{eq:Ageneral}
\end{align}
Thus its symmetric form is
\begin{equation}\label{eq:Qgeneral}
 Q_{\omega,L}=Q^\gamma_{\omega,L}
 -\sum_{\log n<L}\frac{\Lambda(n)}{\sqrt n}
 \cosh(\omega\log n)(T_{\log n}+T_{\log n}^*).
\end{equation}
\end{proposition}
\begin{proof}
The commuting nilpotent delay algebra gives
\[
 B_{\omega,L}=\prod_{\log p<L}
 (I-p^{-1/2-\omega}T_{\log p})
 (I-p^{-1/2+\omega}T_{\log p})^{-1}.
\]
Each inverse is a terminating geometric series. Expanding the logarithm
of one factor, the coefficient of $T_{k\log p}$ is
$(p^{-k/2+k\omega}-p^{-k/2-k\omega})/k$.
Its derivative is
$2(\log p)p^{-k/2}\cosh(k\omega\log p)$.
Differentiating the product in \eqref{eq:Vdef} and using commutation with
causal gamma convolution proves \eqref{eq:Ageneral}. Symmetrization gives
\eqref{eq:Qgeneral}. Only the bounded arithmetic factor has been inverted;
no bounded inverse of the smoothing transfer is asserted.
\end{proof}

\subsection{A bounded quadratic change in the shift}

Set
\begin{equation}\label{eq:rkernel}
 r(t)=e^{t/2}-\frac{e^{-5t/2}}{1-e^{-2t}},\qquad t>0.
\end{equation}
\begin{lemma}[Bounded shift difference]\label{lem:change}
The forms $Q_{\omega,L}$, $0\le\omega\le1/2$, have the common domain
$\DD_{\log,L}$. Their gamma difference has symmetric kernel
\[
 (\cosh(\omega|x-y|)-1)r(|x-y|).
\]
In particular,
\begin{equation}\label{eq:changebound}
 \norm{Q_{\omega,L}-Q_{0,L}}\le C_L\omega^2,
\end{equation}
as a bound on the bounded form difference, where
\begin{align}
 C_L={}&\cosh(L/2)\left(e^{L/2}\frac{L^3}{3}+\frac{L^2}{4}\right)\notag\\
 &+\frac12\sum_{\log n<L}\frac{\Lambda(n)}{\sqrt n}
 (\log n)^2\cosh(\log n/2)\norm{T_{\log n}+T_{\log n}^*}.
 \label{eq:CL}
\end{align}
\end{lemma}
\begin{proof}
In the digamma integral \cite[\S5.9]{DLMF}, form
$a^\gamma_\omega-a^\gamma_0$ before separating the divergent terms. Its
digamma part has inverse Laplace kernel
$-2e^{-t/2}(\cosh(\omega t)-1)/(1-e^{-2t})$.
The rational terms contribute
$4\cosh(t/2)(\cosh(\omega t)-1)$.
Since
$2\cosh(t/2)-e^{-t/2}/(1-e^{-2t})=r(t)$,
symmetrization gives the asserted kernel. It is $O(|x-y|)$ at the diagonal
and extends continuously in $\omega$ through $1/2$. This physical-space
identity avoids taking an incomplete pointwise digamma limit at the endpoint.

Now $|r(t)|\le e^{t/2}+1/(2t)$ and
$0\le\cosh(\omega t)-1\le\omega^2t^2\cosh(t/2)/2$.
The row-integral bound for the symmetric kernel is at most
$\omega^2\cosh(L/2)\int_0^L(t^2e^{L/2}+t/2)\,dt$.
Bounding each arithmetic difference in \eqref{eq:Qgeneral} in the same
way gives \eqref{eq:CL}. Bounded perturbations preserve the closed-form
domain.
\end{proof}

\begin{proposition}[Central coercivity implies contraction]\label{prop:continuation}
If $Q_{0,L}\succeq mI$, then
\begin{equation}\label{eq:continuation}
 \norm{V_{\omega,L}}\le
 \exp(-m\omega+C_L\omega^3/3),\qquad 0<\omega\le1/2.
\end{equation}
\end{proposition}
\begin{proof}
We give the regularity argument because the generator is unbounded and
$V_{0,L}=I$ is a strong, not a norm, limit. Fix a Laplace line $\Re p=\eta>1$.
Uniform gamma asymptotics, together with absolutely convergent zeta Dirichlet
series on this line, give for $0\le s\le1/2$
\[
 |K_s(\eta+i\tau)|\le C(2+|\tau|)^{-s}.
\]
On each $\varepsilon\le s\le1/2$, its first two shift derivatives are
bounded by $C_\varepsilon(1+\log(2+|\tau|))^j(2+|\tau|)^{-\varepsilon}$,
$j=1,2$. Weighted half-line convolution therefore gives norm
differentiability on positive-shift intervals and
\eqref{eq:derivative} after causal restriction. Dominated convergence on
the same line gives $V_{s,L}f\to f$ in $L^2(0,L)$ as $s\downarrow0$.

For $0<r<\min(s,1/2)$ the multiplier maps into $H^r$ locally; multiplication
by an interval indicator is bounded on $H^r(\RR)$ for $r<1/2$
\cite{Strichartz1967}. After restriction and zero extension, the output lies
in $\DD_{\log,L}$. The same estimates with a logarithmic factor justify the
causal generator action and extension of its symmetric form identity from
smooth data. Hence, for every input $f$, on a compact positive-shift interval,
\[
 \frac{d}{ds}\norm{V_{s,L}f}^2
 =-2Q_{s,L}[V_{s,L}f]
 \le-2(m-C_Ls^2)\norm{V_{s,L}f}^2.
\]
Gronwall's inequality, first from $\varepsilon$ to $\omega$ and then with
$\varepsilon\downarrow0$, proves \eqref{eq:continuation}. No norm derivative
at zero or bounded transfer inverse is used.
\end{proof}

\section{An analytic lower bound on the infinite Legendre tail}\label{sec:tail}

Let $\phi_n(x)=\sqrt{(2n+1)/L}\,P_n(2x/L-1)$ and let $P_N$ retain
$\phi_0,\ldots,\phi_{N-1}$. Define causal fractional integration by
\[
 I_\nu f(x)=\frac1{\Gamma(\nu)}\int_0^x(x-y)^{\nu-1}f(y)\,dy.
\]
The following estimate applies to all combinations in the infinite tail.

\begin{lemma}[Fractional tail estimate]\label{lem:fractional}
For $N\ge1$ and $0<\nu<1$,
\begin{equation}\label{eq:fractail}
 \norm{I_\nu(I-P_N)}\le
 (L/2)^\nu\sqrt{\frac{\Gamma(N+1-\nu)}{\Gamma(N+1+\nu)}}.
\end{equation}
Consequently
\begin{equation}\label{eq:I1tail}
 \norm{I_1(I-P_N)}\le\frac{L}{2\sqrt{N(N+1)}}.
\end{equation}
\end{lemma}
\begin{proof}
Termwise beta integration of the terminating Legendre polynomial gives
\[
 I_\nu P_n(2x/L-1)=
 \frac{n!}{\Gamma(n+1+\nu)}x^\nu
 P_n^{(-\nu,\nu)}(2x/L-1).
\]
The polynomial and orthogonality identities may also be read from
\cite[\S\S18.3,18.5,18.17]{DLMF}. For a finite tail combination, use
$x^{2\nu}\le(L/2)^{2\nu}x^\nu(L-x)^{-\nu}$ and Jacobi orthogonality.
Indeed $x(L-x)\le L^2/4$, and the required Jacobi norm is
\[
 \int_0^L x^\nu(L-x)^{-\nu}
 \bigl[P_n^{(-\nu,\nu)}(2x/L-1)\bigr]^2\,dx
 =\frac{L\,\Gamma(n+1-\nu)\Gamma(n+1+\nu)}
 {(2n+1)(n!)^2}.
\]
The squared norm becomes bounded by
\[
 (L/2)^{2\nu}\sum_{n\ge N}|c_n|^2
 \frac{\Gamma(n+1-\nu)}{\Gamma(n+1+\nu)}.
\]
Successive gamma ratios decrease by the factor
$(n+1-\nu)/(n+1+\nu)<1$. Density and boundedness of $I_\nu$ prove
\eqref{eq:fractail}. Letting $\nu\uparrow1$, using convergence of the
fractional kernels in $L^1(0,L)$, gives \eqref{eq:I1tail}.
\end{proof}

The central gamma derivative $U^\gamma=\partial_\omega V^\gamma_{\omega,L}|_0$
has the decomposition
\begin{equation}\label{eq:Udecomp}
 U^\gamma=J_L+S_w,\qquad
 J_L=\left.\partial_\nu[(2\pi)^\nu I_\nu]\right|_0,
 \qquad w(t)=\frac{G_0(t)-1}{t},
\end{equation}
where $S_w$ is causal convolution and
\begin{equation}\label{eq:G0}
 G_0(t)=\frac{e^{t/2}t}{\sinh t}-4t\cosh(t/2)
 =\sum_{j\ge0}g_jt^j.
\end{equation}
To verify \eqref{eq:Udecomp}, the inverse kernel of $-a^\gamma_0$, away
from the origin, is
$e^{t/2}/\sinh t-4\cosh(t/2)$; subtracting the finite-part kernel of
$J_L$ leaves $w$. The constant term agrees with the Laplace asymptotic
$-a^\gamma_0(p)=\log(2\pi/p)+O(p^{-1})$.
In particular $g_0=1$, $g_1=-7/2$, and $g_2=-1/24$.

\begin{proposition}[Central infinite-tail floor]\label{prop:tail}
Suppose
$K_L\ge |w(0)|+\int_0^L|w'(t)|\,dt$. Then, on
$(I-P_N)\DD_{\log,L}$,
\begin{equation}\label{eq:afloor}
 Q_{0,L}\succeq a_{N,L}I,
\end{equation}
where
\begin{align}
 a_{N,L}={}&H_N-\gamma-\log(\pi L)
 -\frac{K_LL}{2\sqrt{N(N+1)}}\notag\\
 &-\sum_{\log n<L}\frac{\Lambda(n)}{\sqrt n}
 \norm{T_{\log n}+T_{\log n}^*}.
 \label{eq:tailformula}
\end{align}
\end{proposition}
\begin{proof}
Apply \eqref{eq:fractail} to a polynomial tail vector, multiply by
$(2\pi)^\nu$, square, and differentiate at $\nu=0$. Both sides have the
same value there. The result is
$-\Rea\ip f{J_Lf}\ge(H_N-\gamma-\log(\pi L))\norm f^2$.
The differentiation is legitimate on each polynomial, whose derivative
contains only polynomials and logarithmic endpoint terms in $L^2$.

Next $S_w=(w(0)I+S_{w'})I_1$. Young's inequality and \eqref{eq:I1tail}
give
\[
 \norm{S_w(I-P_N)}\le\frac{K_LL}{2\sqrt{N(N+1)}}.
\]
Since $Q^\gamma_0=-\Rea U^\gamma$ in form sense, this is a lower-bound
loss. Bound the remaining symmetric delays individually in norm.
Lemma~\ref{lem:core} extends the resulting polynomial inequality to the
complete tail form domain.
\end{proof}

For clarity, the delay norm in \eqref{eq:tailformula} is not automatically 2.
Decompose $x$ into $x=r+kd$, $0<r<d$. On each fiber,
$T_d+T_d^*$ is the adjacency matrix of a finite path. If
$q=\lceil L/d\rceil$, its norm is
\begin{equation}\label{eq:chainnorm}
 \norm{T_d+T_d^*}=2\cos\frac{\pi}{q+1}
 \quad(d<L),
\end{equation}
with threshold equality using the shorter path, since the exceptional fiber
has measure zero. In particular the norm is 1 for $d<L\le2d$ and
$\sqrt2$ for $2d<L\le3d$. The reviewed implementation evaluates the
cosine formula for every finite path length. Logarithmic endpoints use
exact integer-power comparisons; unresolved ball comparisons at other
horizons stop the calculation.

\subsection{A uniform analytic profile remainder}

\begin{lemma}[Profile enclosure]\label{lem:profile}
For $0<L<3$ and an integer $M\ge1$, let $\widetilde U$ retain the terms
$g_j$, $j\le M$, in \eqref{eq:Udecomp}, leaving $J_L$ and the arithmetic
delays exact. Then the bounded difference satisfies
\begin{equation}\label{eq:eta}
 \norm{U-\widetilde U}\le\eta_M,
 \qquad \eta_M=\frac{256(L/3)^{M+1}}{(M+1)(1-L/3)}.
\end{equation}
One may use
\begin{equation}\label{eq:Kbound}
 K_L=\sum_{j=1}^M|g_j|L^{j-1}
 +\frac{256}{3}\frac{(L/3)^M}{1-L/3}
\end{equation}
in Proposition~\ref{prop:tail}.
\end{lemma}
\begin{proof}
On $|z|=3$, the Euler product for $\sinh z/z$
\cite[\S4.36]{DLMF} gives $|z/\sinh z|\le3/\sin3<24$.
Since $e^{3/2}<5$, \eqref{eq:G0} implies
$|G_0(z)|<5(24+12)<256$. Thus $|g_j|\le256\,3^{-j}$.
The $L^1$ norm of the omitted convolution kernel is bounded by
$\sum_{j>M}|g_j|L^j/j\le\eta_M$, proving \eqref{eq:eta}.
For $w$, the constant term contributes $|g_1|$, while integration of each
derivative term gives $|g_j|L^{j-1}$. Summing the omitted geometric series
proves \eqref{eq:Kbound}.
\end{proof}

\section{Exact reflection-parity coupling and the Schur test}\label{sec:schur}

Write $P=P_N$, let $P_r$ be its parity-$r$ part ($r=\pm1$), and denote
the central model form by $\widetilde Q=-(\widetilde U+\widetilde U^*)/2$.
The polynomial space lies in the operator domain: the fractional derivative
of a polynomial has only square-integrable logarithmic endpoint terms,
and all other terms are bounded convolutions or delayed polynomials.
Thus $Q_0P$ and $\widetilde QP$ are bounded finite-column operators.

Define the retained matrix and two full-output Grams by
\begin{equation}\label{eq:grams}
 q=P\widetilde QP,\qquad
 F=(\widetilde UP)^*(\widetilde UP),\qquad
 C=(\widetilde UP)^*R_L(\widetilde UP).
\end{equation}
In this notation, $F$ and $C$ are finite matrices but the output in their
definition has \emph{not} been projected onto $P\HH_L$.

\begin{lemma}[Exact central-parity leakage]\label{lem:parity}
If $q_r,F_r,C_r$ denote the same-parity blocks in \eqref{eq:grams}, then
\begin{equation}\label{eq:leakage}
 E_r=\frac12(F_r+rC_r)-q_r^2
 =\bigl((I-P)\widetilde QP_r\bigr)^*
   \bigl((I-P)\widetilde QP_r\bigr).
\end{equation}
\end{lemma}
\begin{proof}
Causal convolution gives $\widetilde U^*=R_L\widetilde UR_L$.
For $R_Lf=rf$,
$\widetilde Qf=-(\widetilde Uf+rR_L\widetilde Uf)/2$.
Taking the full-output squared norm yields
$\norm{\widetilde Qf}^2=\ip f{(F_r+rC_r)f}/2$.
Subtract the retained-output norm $\norm{q_rf}^2$ and polarize.
\end{proof}

\begin{proposition}[Certified finite test]\label{prop:schur}
Suppose the exact central tail is at least $aI$, $a>m$, and
$\norm{\widetilde UP}<K$. If both finite inequalities
\begin{equation}\label{eq:finite-test}
 (a-m)(q_r-mI)-E_r-\epsilon I\succeq0,
 \qquad r=\pm1,
\end{equation}
hold, where
\begin{equation}\label{eq:eps}
 \epsilon=(|a-m|+4K)\eta_M+2\eta_M^2,
\end{equation}
then $Q_{0,L}\succeq mI$ on the full form domain.
\end{proposition}
\begin{proof}
The profile difference has norm at most $\eta_M$, as does its symmetric
part. On each input parity, $\widetilde QP_r$ is an output-parity projection
of $-\widetilde UP_r$, and hence has norm below $K$. The full squared-output
Gram and the squared retained matrix each change by at most
$2K\eta_M+\eta_M^2$. The linear head term changes by at most
$|a-m|\eta_M$. Therefore \eqref{eq:finite-test} implies
\[
 (a-m)(Q_{P,r}-mI)-B_r^*B_r\succeq0,
 \quad B_r=(I-P)Q_0P_r,
\]
for the exact central form.

For an arbitrary input in one sector, write it as $h+t$ with $h$ retained
and $t$ in the form-domain tail. The quadratic form of $Q_0-mI$ is bounded
below by
\[
 \ip h{(Q_{P,r}-mI)h}
 +2\Rea\ip{B_rh}t+(a-m)\norm t^2.
\]
Completing the square, or minimizing over $t$, leaves
$\ip h{(Q_{P,r}-mI-(a-m)^{-1}B_r^*B_r)h}\ge0$.
The two orthogonal reflection sectors exhaust the complex Hilbert space.
\end{proof}

The use of the exact parity leakage is consequential. Replacing it by
the full leakage of $\widetilde U$ also bounds the coupling, but charges
irrelevant output parity against the available tail margin. Equation
\eqref{eq:leakage} removes that loss without discarding any coupling of
the central form.

\begin{remark}[A smaller error for future calculations]\label{rem:sharper-error}
The certificates below keep the conservative error \eqref{eq:eps}.
One can instead set
\[
 \epsilon_r=(|a-m|+2\kappa_r)\eta_M+\eta_M^2,
 \qquad \kappa_r\ge\norm{(I-P)\widetilde QP_r}.
\]
Indeed the omitted-output column operator changes in norm by at most
$\eta_M$, so its squared Gram changes by at most
$2\kappa_r\eta_M+\eta_M^2$ directly. There is no need to bound the full
Gram and the squared head separately. A certified upper bound on
$\sqrt{\tr E_r}$ supplies $\kappa_r$, since $E_r$ is an exact positive
semidefinite Gram. Any numerical upper bound used here must include the
ball uncertainty; this refinement is not used to justify the reported runs.
\end{remark}

\section{Finite construction and the two certificates}\label{sec:computation}

\subsection{Polynomial and logarithmic output formulas}

Use the unitary coordinate $u=x/L$, so the normalized basis is
$\sqrt{2n+1}\,p_n(u)$, where
\begin{equation}\label{eq:legpoly}
 p_n(u)=\sum_{k=0}^n(-1)^{n+k}\binom nk\binom{n+k}k u^k.
\end{equation}
With $\ell_L=\gamma+\log(2\pi L)$, the gamma derivative of a monomial is
\begin{equation}\label{eq:monomial}
 U^\gamma u^k=u^k(\log u+\ell_L-H_k)
 +\sum_{j\ge1}g_jL^j\frac{(j-1)!k!}{(j+k)!}u^{j+k}.
\end{equation}
This follows by differentiating the beta integral for $I_\nu u^k$ and
integrating each term of $S_w$. For the model, the sum stops at $j=M$.
Every output therefore has the form
\begin{equation}\label{eq:output}
 \widetilde Up_n(u)=p_n(u)\log u+q_n(u)
 +\sum_{\log j<L}\frac{2\Lambda(j)}{\sqrt j}
   \mathbf1_{u>\delta_j}p_n(u-\delta_j),
 \quad \delta_j=\frac{\log j}{L}.
\end{equation}
The polynomial $q_n$ is supplied by \eqref{eq:monomial}.
The head, $F$, and $C$ are integrated using exact formulas for moments of
$1,\log u,(\log u)^2,\log(1-u)$, and $\log u\log(1-u)$, including
truncated integration intervals. Appendix~\ref{app:moments} lists them.
There is no quadrature error and no finite-output truncation.

For the reflected delayed products, the nonempty integration window is
$(\delta_i,1-\delta_j)$, equivalently $\log(ij)<L$.
At $L=9/5$ the active generator indices are $2,3,4,5$, and the ordered
overlap pairs are exactly $(2,2),(2,3),(3,2)$.

\subsection{Enclosed arithmetic at \texorpdfstring{$9/5$}{9/5}}

The standalone implementation uses Arb real balls \cite{Johansson2017}
through python-flint \cite{PythonFlint}. Integer and rational coefficients
in \eqref{eq:G0}, \eqref{eq:legpoly}, and \eqref{eq:monomial} are formed
exactly. All subsequent conversions, elementary constants, logarithms,
square roots, moments, polynomial reflections, and matrix operations are
performed with balls at 1536-bit working precision. Exact symmetries are
imposed only after checking consistency of the enclosures.

The matrix $F$ encloses the model-output Gram, and its trace is checked
below $10^7$. Hence $\norm{\widetilde UP}<10000$, allowing $K=10000$
in \eqref{eq:eps}. The computed trace is approximately $1743.1632085$;
the deliberately loose bound avoids making this part of the proof delicate.
An outward lower value of \eqref{eq:tailformula} and an outward upper value
of \eqref{eq:eta} are used. An outward upper value of \eqref{eq:eps} is
then subtracted from each model Schur matrix.

For a symmetric ball matrix $M$, the sign test uses the recurrence
\begin{align*}
 d_i&=M_{ii}-\sum_{k<i}\ell_{ik}^2d_k,\\
 \ell_{ji}&=\frac{M_{ji}-\sum_{k<i}\ell_{jk}\ell_{ik}d_k}{d_i}\quad(j>i).
\end{align*}
Every pivot must be strictly positive as a ball. All divisions then have
denominators excluding zero. Induction encloses the exact $LDL^*$ factors
of the true symmetric matrix; positive pivots imply its positivity.
Ordinary floating-point eigenvalues never enter this decision.

\begin{theorem}[Computer-assisted matrix inequality]\label{thm:computation}
At $L=9/5$, $N=128$, $M=180$, and $m=10^{-26}$, both inequalities
\eqref{eq:finite-test} hold with $K=10000$, the tail floor
\eqref{eq:tailformula}, and the error \eqref{eq:eps}.
In particular, $Q_{0,9/5}\succeq10^{-26}I$.
\end{theorem}
\begin{proof}[Computer-assisted proof]
The supplied builder evaluates \eqref{eq:output}, the moment formulas in
Appendix~\ref{app:moments}, and the bounds of Section~\ref{sec:tail} with
ball arithmetic. The following rationally rounded inequalities are checked:
\begin{align*}
 a_{128,9/5}&>0.69124708426711,&K_{9/5}&<5.482684,\\
 \eta_{180}&<2.477\times10^{-40},&
 \epsilon&<9.908\times10^{-36}.
\end{align*}
The arithmetic norm term in \eqref{eq:tailformula} is approximately
$2.39376738699$. Both $64\times64$ ball LDL tests pass. Their minimum
pivots are displayed in Table~\ref{tab:certificate}. Applying
Proposition~\ref{prop:schur} gives the full-operator conclusion.
\end{proof}

\begin{table}[ht]
\centering
\caption{Certificate data. Displayed pivot values are rounded summaries,
not spectral lower bounds. The sign decisions use complete balls.}
\label{tab:certificate}
\begin{tabular}{@{}lrr@{}}
\toprule
Quantity & Even sector & Odd sector\\
\midrule
Retained dimension & 64 & 64\\
Target central floor & $10^{-26}$ & $10^{-26}$\\
Minimum LDL pivot & $5.6796287856\times10^{-5}$ & $1.8170726676\times10^{-4}$\\
Every pivot strictly positive & Yes & Yes\\
Replay from saved balls & Pass & Pass\\
\bottomrule
\end{tabular}
\end{table}

The largest radius among the four saved matrix enclosures was less than
$9.341\times10^{-215}$. This finite-arithmetic uncertainty is propagated
through the sign tests; it is separate from the analytic remainder.
The matrix archive stores dyadic midpoint/radius pairs and checks that every
reconstructed ball contains the original. Reconstructing a radius may enlarge
it outward. A separate process loaded this archive and repeated both LDL
tests successfully.

\begin{proof}[Proof of Theorem~\ref{thm:main}]
Theorem~\ref{thm:computation} establishes \eqref{eq:maincentral} at $L=9/5$.
Zero extension into a larger interval preserves the central form: the
Fourier representation and its translation-invariant kernels, including
the delay correlations, do not depend on the chosen containing horizon.
Thus the bound holds at every shorter horizon.

Evaluation of \eqref{eq:CL} at $9/5$ gives $C_{9/5}<10.298615<11$.
The exact rational inequality
\[
 10^{-26}-\frac{11}{3}(3\times10^{-14})^2
 =6.7\times10^{-27}>5\times10^{-27}
\]
and Proposition~\ref{prop:continuation} prove \eqref{eq:maintransfer}
at the outer horizon. Causality makes every shorter transfer a compression,
so it has no larger norm. Self-adjointness of $H=VR_L$ and
$\norm H=\norm V$ give \eqref{eq:mainreflection}.
\end{proof}

\subsection{What the larger head resolved}

The preceding 96-mode experiment at the same horizon did not pass the
sufficient Schur test. To explain the improvement, Table~\ref{tab:diagnostic}
compares the scalar tail floor with the diagnostic largest eigenvalue of
$q_r^{-1/2}E_rq_r^{-1/2}$. The latter measures the tail floor required by
the model test at level zero. It was computed using a high-precision
congruence before conversion to floating point and is a diagnostic only.

\begin{table}[ht]
\centering
\caption{Head/tail comparison at $L=9/5$. Only the final column represents
the complete certified sign test at $m=10^{-26}$.}
\label{tab:diagnostic}
\begin{tabular}{@{}rrrrl@{}}
\toprule
$N$ & Tail floor & Required, even & Required, odd & Complete test\\
\midrule
96 & 0.39212903 & 0.55845 & 0.54762 & Not certified\\
128 & 0.69124708 & 0.49054476 & 0.47957150 & Pass\\
\bottomrule
\end{tabular}
\end{table}

Thus increasing the retained space improved both the tail floor and the
remaining coupling. Failure of the earlier sufficient inequality was not
evidence of a negative direction of the full form. The old matrix construction
used a separate aggregate arithmetic budget restricted to at most 96 modes;
that budget was not extended here. The new direct-ball calculation is
independent of it. As a normalization diagnostic, all shared entries of the
new head and full-output Grams agree with the earlier 96-mode point matrices
to within $3\times10^{-217}$ after conversion of the causal/Hankel signs.
Those old entries are not inputs to the new positivity proof.

\subsection{Extension to the exact endpoint \texorpdfstring{$\log7$}{log 7}}

\begin{theorem}[Computer-assisted endpoint inequality]\label{thm:log7}
At $L=\log7$, $N=128$, $M=220$, and $m=10^{-34}$, both inequalities
\eqref{eq:finite-test} hold with $K=10000$ and the same analytic reduction.
Consequently $Q_{0,\log7}\succeq10^{-34}I$.
\end{theorem}
\begin{proof}[Computer-assisted proof]
The direct-Arb construction at 1792-bit precision verifies
\begin{align*}
 a_{128,\log7}&>0.60748624181467,& K_{\log7}&<5.839924,\\
 \eta_{220}&<9.342\times10^{-42},&
 \epsilon&<3.737\times10^{-37}.
\end{align*}
The full Gram trace is less than $1712$, hence less than $10^7$.
Both 64-dimensional ball LDL tests pass. The minimum pivots are approximately
$1.6916652823\times10^{-5}$ (even) and $9.0360554449\times10^{-5}$ (odd).
These pivots are sign-test diagnostics, not central spectral margins.
At the exact endpoint $T_{\log7}=0$ almost everywhere. The active
generator indices remain $2,3,4,5$, and the reflected ordered overlaps
remain $(2,2),(2,3),(3,2)$. Integer comparisons implement endpoint equality.
Proposition~\ref{prop:schur} completes the full-operator argument.
\end{proof}

\begin{proof}[Proof of Theorem~\ref{thm:endpoint-main}]
Theorem~\ref{thm:log7} and zero extension give the central bound.
Direct ball evaluation of \eqref{eq:CL} gives
$C_{\log7}<13.541169<14$. The exact rational calculation
\[
 10^{-34}-\frac{14}{3}(3\times10^{-18})^2
 =5.8\times10^{-35}>5\times10^{-35}
\]
and Proposition~\ref{prop:continuation} give the transfer bound.
Reflection and compression give the remaining conclusions.
\end{proof}

The 128-mode tail floor is close to the required model floor at this
endpoint: the even and odd relative-coupling diagnostics are approximately
$0.59479118$ and $0.58991735$, compared with $a\approx0.60748624$.
Their ratios to $a$ are approximately $0.97910$ and $0.97108$.
Thus the suggested 128-mode calculation succeeds without increasing the
retained dimension, but leaves less room in this sufficient test than at
$9/5$.

\subsection{Independent reconstruction during review}

A second implementation, \path{independent_arb.py}, was begun from the
printed equations before the supplied numerical source became available.
It normalizes the polynomial rows before assembly and computes reflected
gamma--delay integrals over $(0,1-\delta)$, rather than using the supplied
builder's tail moments over $(\delta,1)$. It reproduces both full-operator
tests and their displayed pivots, using Python 3.10.0, python-flint 0.9.0,
and FLINT 3.6.0. All entries of the retained head and both full-output Grams
have overlapping enclosures between implementations at both horizons.
Replaying the independently saved Schur matrices also passes.
This is an independent implementation check of the same analytic reduction;
both programs use Arb and neither constitutes independent human review or
formal verification.

\section{Storage, scope, and the next depth problem}\label{sec:storage}

\begin{corollary}[Cumulative storage]\label{cor:storage}
For any spatial split of $(0,L)$ with $L\le9/5$, write
\[
 V_{\omega,L}=\begin{pmatrix}X&0\\Y&Z\end{pmatrix}.
\]
For $0<\omega\le3\times10^{-14}$,
\begin{equation}\label{eq:storage}
 I-X^*X-Y^*(I-ZZ^*)^{-1}Y
 \succeq(1-e^{-10^{-26}\omega})I.
\end{equation}
\end{corollary}
\begin{proof}
Let $\rho=e^{-5\times10^{-27}\omega}$. Theorem~\ref{thm:main} gives
$I-V^*V\succeq(1-\rho^2)I$, and $\norm Z\le\rho<1$.
The Schur complement of the lower-right block $I-Z^*Z$ is
\[
 I-X^*X-Y^*Y-Y^*Z(I-Z^*Z)^{-1}Z^*Y
 =I-X^*X-Y^*(I-ZZ^*)^{-1}Y.
\]
For a fixed first-block vector, minimizing the defect quadratic form over
the second-block vector leaves a value at least
$(1-\rho^2)$ times the first-block norm squared. This proves
\eqref{eq:storage}.
\end{proof}

In particular, the split $9/5=\log5+(9/5-\log5)$ includes the first
mixed-prime region. The corollary is obtained by certifying the whole larger
operator first. It does not show that positivity of a previous slab alone
forces the next storage update to be positive. That distinction separates
the present finite-depth result from the intended induction problem.

The same proof and Theorem~\ref{thm:endpoint-main} give, for $L\le\log7$
and $0<\omega\le3\times10^{-18}$,
\[
 I-X^*X-Y^*(I-ZZ^*)^{-1}Y\succeq(1-e^{-10^{-34}\omega})I.
\]

\subsection{A diagonal formulation of the all-depth target}

\begin{proposition}\label{prop:diagonal}
Suppose there exist $L_j\to\infty$ and $\omega_j\downarrow0$ such that
$\norm{V_{\omega_j,L_j}}\le1$ for every $j$. Then the central Weil form
is nonnegative on every compactly supported smooth test function.
By Weil's criterion, such a sequence would imply RH.
\end{proposition}
\begin{proof}
Fix a finite horizon $L$ and a smooth test $f$ compactly supported inside it.
For all sufficiently large $j$, causal compression gives
$\norm{V_{\omega_j,L}f}\le\norm f$.
The scalar derivative at zero on this smooth core follows by dominated
convergence on a far right Laplace line: the transform of the input is
rapidly decreasing, while the differentiated multiplier grows only
logarithmically. It is
\[
 \lim_{s\downarrow0}\frac{\norm f^2-\norm{V_{s,L}f}^2}{2s}
 =Q_{0,L}[f].
\]
The quotients along $s=\omega_j$ are nonnegative, so their limit is
nonnegative. The horizon and compact test were arbitrary; the final
implication is the classical Weil criterion
\cite{ConnesConsani2020,Conrey2019}.
\end{proof}

The proposition requires neither a depth-independent coercivity margin nor
a fixed positive shift interval. Nevertheless, constructing the entire
sequence remains the missing all-depth statement; the single horizon
certified here is not such a sequence.

\subsection{An exact storage update with memory}\label{sec:recursion}

Fix a positive shift and suppose $D_X=I-X^*X\succ0$ and $\norm Z<1$,
with bounded positive inverses. Put
\[
 W=I-Z^*Z,\qquad
 \mathcal C=(I-ZZ^*)^{-1/2}YD_X^{-1/2},\qquad
 S=D_X^{1/2}(I-\mathcal C^*\mathcal C)D_X^{1/2}.
\]
For $V=\left(\begin{smallmatrix}X&0\\Y&Z\end{smallmatrix}\right)$ one has
the exact factorization
\begin{equation}\label{eq:storage-factor}
 I-V^*V=
 \begin{pmatrix}I&-Y^*ZW^{-1}\\0&I\end{pmatrix}
 \begin{pmatrix}S&0\\0&W\end{pmatrix}
 \begin{pmatrix}I&0\\-W^{-1}Z^*Y&I\end{pmatrix}.
\end{equation}
Multiplication verifies the identity, using
$I+Z(I-Z^*Z)^{-1}Z^*=(I-ZZ^*)^{-1}$.
Thus contraction is equivalent to $\norm{\mathcal C}\le1$.
If $\norm{\mathcal C}<1$, the two triangular factors and the storage
factor provide the complete inverse defect for the next append operation.
If only equality is certified, the defect may be singular and this
bounded-inverse recursion cannot automatically continue.

This supplies an algebraic recursion at arbitrary finite depth. Its missing
arithmetic step is a bound on $\mathcal C$ that propagates for the actual
prime--gamma kernel. A scalar lower bound on $D_X$ discards orientation:
for $Z=0$ and $Y=(c,0)$, the defects
$\operatorname{diag}(\varepsilon,1)$ and
$\operatorname{diag}(1,\varepsilon)$ have the same least eigenvalue, but
their updates are positive precisely when $c^2\le\varepsilon$ and
$c^2\le1$, respectively. Each defect is realized by a diagonal contraction
$X$. Taking $\varepsilon=1/4$ and $c^2=1/2$ gives a failing and a passing
update with identical old scalar margins. A useful induction therefore
needs the defect factor and its coupling to the next output.

\subsection{Nested Schur elimination for a structured tail}

There is a parallel recursion in retained degree. Split a form into an old
head, a finite buffer, and the remaining tail, and write
\[
 Q_0-mI=\begin{pmatrix}
 A&B^*&C^*\\B&D&E^*\\C&E&F
 \end{pmatrix},\qquad F\succeq\alpha I,\quad\alpha>0.
\]
The finite head and buffer lie in the operator domain, so the displayed
columns into the tail are bounded. Replacing $F$ by $\alpha I$ gives a
lower form bound. Eliminating that tail yields
\[
 \widehat A=A-C^*C/\alpha,\quad
 \widehat B=B-E^*C/\alpha,\quad
 \widehat D=D-E^*E/\alpha.
\]
Consequently the two tests
\begin{equation}\label{eq:nested-schur}
 \widehat D\succ0,\qquad
 \widehat A-\widehat B^*\widehat D^{-1}\widehat B\succeq0
\end{equation}
are sufficient for $Q_0\succeq mI$. The proof is two successive
completions of the square. The cross-Gram $C^*E$ must be retained;
separate norm bounds on $C$ and $E$ discard its cancellation.
Full-output parity Grams for the enlarged retained space supply all three
products $C^*C,C^*E,E^*E$. The same construction works in each parity sector.
With the same enlarged space and scalar far-tail floor, this is a reordered
version of the existing test, not an intrinsically stronger theorem. It
organizes factor reuse and identifies where an improved operator lower bound
for $F$ would strengthen the result. Such a bound requires its own proof.

\subsection{Limitations and further calculations}

The scalar tail bound pays separately for every arithmetic delay norm.
This is sufficient through $9/5$ once the retained space is enlarged, but
does not explain the prime--gamma cancellation at unbounded depth. The
small-shift interval in Theorem~\ref{thm:main} is a consequence of the stated
coercivity floor and a quadratic perturbation estimate; it is not asserted
to be sharp. It does not connect to the previously investigated interior
shift ranges of the second slab.

The $\log7$ endpoint target is now certified. A next test is $L=\log8$,
which activates the $n=7$ prime while keeping the $n=8$ endpoint delay zero
almost everywhere. The active generator indices would be $2,3,4,5,7$.
No certificate at that horizon is asserted here. Choose the profile degree
and working precision from its own error budget before increasing the head.
The smaller error in Remark~\ref{rem:sharper-error} and the buffer
formulation \eqref{eq:nested-schur} are concrete candidates to test.

For horizons beyond the radius of the present origin-centered profile,
one can seek validated piecewise polynomial approximations to $w$ on the
positive real axis, retaining the logarithmic kernel exactly. The complex
poles of $G_0$ lie at nonzero integer multiples of $i\pi$, whereas $w$ is
smooth at every nonnegative real argument. Taylor expansions about positive
centers or Chebyshev pieces can therefore cover any fixed finite interval,
provided their errors and derivative variations are enclosed. Their convolution
outputs introduce additional piecewise polynomial terms and integration
breakpoints, which must all be retained in the ordinary and reflected Grams.
This removes the specific $L<3$ profile restriction in principle; it does not
prove coercivity at larger depth or bound the cost of doing so.

\appendix
\section{Moment formulas and assembly conventions}\label{app:moments}

The code constructs all polynomial coefficients and integrates them against
the moments below. For an integer $k\ge0$, put $s=k+1$,
$H_s=\sum_{j=1}^s1/j$, and $H_s^{(2)}=\sum_{j=1}^s1/j^2$.
The full-interval formulas are
\begin{align}
 \int_0^1u^k\,du&=\frac1s,&
 \int_0^1u^k\log u\,du&=-\frac1{s^2},\notag\\
 \int_0^1u^k(\log u)^2\,du&=\frac2{s^3},&
 \int_0^1u^k\log(1-u)\,du&=-\frac{H_s}{s},\label{eq:moments}\\
 \int_0^1u^k\log u\log(1-u)\,du
 &=\frac{H_s}{s^2}-\frac{\pi^2/6-H_s^{(2)}}s.\notag
\end{align}
They follow by differentiating the beta integral; the first three also
follow directly by differentiating $1/s$.
For $0<\delta<1$,
\begin{align}
 \int_\delta^1u^k\,du&=\frac{1-\delta^s}s,\notag\\
 \int_\delta^1u^k\log u\,du
 &=-\frac{1-\delta^s}{s^2}-\frac{\delta^s\log\delta}s,\label{eq:tailmoments}\\
 \int_\delta^1u^k\log(1-u)\,du
 &=\frac{(1-\delta^s)\log(1-\delta)-H_s+\sum_{j=1}^s\delta^j/j}{s}.
 \notag
\end{align}
Integration by parts gives the last formula. Products of delayed
polynomials use
$\int_\delta^b u^k\,du=(b^s-\delta^s)/s$ on the appropriate nonempty
window. For reflection, expand $(1-u)^k$ by the exact binomial theorem.

To recover the gamma polynomial $q_n$ in \eqref{eq:output}, multiply each
coefficient of $u^k$ in \eqref{eq:legpoly} by $\ell_L-H_k$, and add
its contributions $g_jL^j(j-1)!k!/(j+k)!$ at degree $k+j$.
The beta factor is computed rationally with the recurrence
$\beta_{1,k}=1/(k+1)$,
$\beta_{j,k}=\beta_{j-1,k}(j-1)/(j+k)$.
The series $g_j$ is obtained by reciprocal power-series inversion of
$\sinh t/t$, multiplication by $e^{t/2}$, and subtraction of
$4t\cosh(t/2)$, all over the rationals.

The unit-coordinate integrals acquire factors $\sqrt{(2i+1)(2j+1)}$ from
normalization. The saved matrices use the causal derivative $U$.
If instead $H_1=UR_L$ is used, the two full-output Grams acquire the
entry sign $(-1)^{i+j}$; this sign is 1 within either same-parity block.
The central head and central leakage are unchanged.

\section{Reproducibility and review map}\label{app:reproduce}

The accompanying review package contains the LaTeX source, bibliography,
compiled manuscript, continuation guide, and the full numerical supplement.
The core numerical files under \path{numerics/} are:
\begin{itemize}
\item \path{certify_arb.py}: exact-rational profiles, full-output ball
  construction, analytic tail bounds, and ball LDL;
\item \path{analyze_certificate.py}: rational continuation checks and
  diagnostic relative coupling;
\item \path{central_matrices.json.gz}: the four saved ball matrices and
  analytic-bound enclosures;
\item \path{central_certificate.json}: both complete pivot lists and
  the matrix-archive SHA256;
\item \path{analysis.json}: continuation validation and separately labelled
  diagnostics.
\item \path{independent_arb.py}: the separately written construction from
  the printed formulas; its saved matrices and pivot records are included.
\item \path{review_checks.py}: moment, arithmetic-generator, parity, and
  failure-control checks, with results under \path{review/}.
\end{itemize}

The recorded numerical environment was Python 3.12.14, python-flint 0.9.0,
FLINT 3.6.0, and NumPy 2.3.5. NumPy is used only for diagnostic eigenvalues.
The builder uses the Python standard library and python-flint. Run from
the numerical supplement directory, with assertions enabled:
\begin{verbatim}
python certify_arb.py --N 128 --M 180 --bits 1536 \
  --horizon 9/5 --floor 1e-26 --output output/rebuilt
python certify_arb.py --reuse --bits 1536 --floor 1e-26 \
  --output output/length_1p8_N128
python analyze_certificate.py output/length_1p8_N128 \
  --shift 3e-14 --decay 5e-27 --generator-bound 11
python certify_arb.py --N 128 --M 220 --bits 1792 \
  --log-horizon 7 --floor 1e-34 --output output/log7_rebuilt
python analyze_certificate.py output/log7_rebuilt \
  --shift 3e-18 --decay 5e-35 --generator-bound 14
\end{verbatim}
A full rebuild checks the construction as well as the final sign tests.
Replay checks the signs from the saved construction enclosures. The latter
does not replace review of the analytic formulas. Do not run Python with
\texttt{-O}. Build times are machine-dependent; the recorded matrix
construction took about 64 seconds, before archive writing and final LDL.
The reviewed commands reject optimized Python mode and return a nonzero
exit status if certification fails. Replay recomputes analytic bounds,
the full-Gram trace check, and central leakage from the saved head and Grams.
The analyzer checks archive hashes and metadata and repeats the sign tests
before deriving continuation. Its shift and decay inputs must be positive;
its checked generator constant is specific to the loaded horizon. An omitted
shift is chosen by an exact rational halving rule, rather than reusing the
$9/5$ constants at other horizons.

The largest mathematical review priorities are the normalization in
\eqref{eq:weil} and \eqref{eq:Ageneral}, the common-domain energy argument,
the operator-level fractional tail estimate, and the full-output parity
identity. The numerical review should independently check the moment
assembly, outward ball serialization, and the subtraction of
\eqref{eq:eps} before each LDL test. No first-slab numerical archive or
external compact-window certificate is required for this proof.

\section*{Acknowledgment of computational assistance}
This draft and its numerical implementation were prepared with substantial
assistance from OpenAI language models. The record distinguishes analytic
proofs, validated arithmetic, and exploratory diagnostics. The manuscript
is being circulated for mathematical and computational review.

\bibliographystyle{plain}
\bibliography{finite_horizon_weil}
\end{document}
